\chapter{Fundamentação Teórica}

\section{Medição de Energia}

Para medir o consumo de energia, há duas metodologias predominantes na
literatura. A primeira consiste no uso de \textit{hardware} suplementar para
medir o consumo diretamente da tomada, resultando em uma medida fiel do consumo
real de energia \cite{energy-java-collections-hasan}. A segunda abordagem
consiste em utilizar estimativas de consumo geradas pela própria máquina. Neste
trabalho, optou-se pela segunda metodologia. A ferramenta utilizada para
estimar o consumo de energia durante a execução dos \textit{benchmarks} foi o
Intel Running Average Power Limit (RAPL) \cite{rapl}. O RAPL é uma
funcionalidade introduzida nas CPUs da Intel a partir da arquitetura Sandy
Bridge. Ele utiliza registradores de 32 \textit{bits} para armazenar a
informação do consumo de energia desde a inicialização da máquina; esses
registradores são atualizados aproximadamente a cada 1 milissegundo. Por meio
dessa ferramenta, é possível obter estimativas do consumo de energia de
diferentes componentes, como o pacote completo da CPU, seus \textit{cores} e a
DRAM. O RAPL apresenta precisão suficiente para servir como uma aproximação
fiel do consumo real de energia: \citeonline{rapl-in-action} encontraram
correlação de 0,99 entre o valor estimado pelo RAPL e o valor real medido
externamente. Além disso, os autores mostram que a sobrecarga introduzida pelo
próprio mecanismo é negligenciável. Outra vantagem do uso do RAPL em relação a
métodos diretos é sua maior granularidade, permitindo focalizar partes
específicas do \textit{hardware}, enquanto medições externas normalmente
observam apenas o sistema como um todo. Sendo assim, o RAPL constitui um método
confiável e conveniente para obter medições de energia.

É possível acessar os registradores do RAPL diretamente para obter os valores
de energia. Entretanto, em vez desse método, foi empregada a ferramenta JRAPL
\cite{jrapl}. O JRAPL é um \textit{wrapper} que permite acessar os
registradores diretamente a partir de código Java. A sobrecarga decorrente
desse acesso é negligenciável, de modo que seu uso não compromete as medições
\cite{programming_lang_energy}. A integração direta com o código Java é uma
vantagem importante do JRAPL em relação ao acesso manual aos registradores, pois
elimina problemas de sincronização entre o código dos \textit{benchmarks} e a
coleta das estimativas, garantindo que as leituras sejam obtidas imediatamente
antes e imediatamente depois das execuções.

\section{Execução de benchmarks na JVM}

A execução de \textit{benchmarks} em Java apresenta desafios adicionais em
comparação a outras linguagens. Isso se deve principalmente ao ambiente de
execução da linguagem, isto é, à Java Virtual Machine (JVM).

Ao contrário de linguagens compiladas estaticamente, o comportamento de uma
aplicação Java pode se alterar ao longo da própria execução. Esse efeito
decorre da compilação \textit{just-in-time} (JIT), do perfilamento adaptativo
do código, do carregamento dinâmico de classes e da atuação do coletor de lixo.
Como consequência, as primeiras iterações de um \textit{benchmark}
frequentemente representam uma fase de aquecimento, e não o comportamento
estável do sistema \cite{vms-hot-and-cold, steady-state-assessment}. Além
disso, mesmo em ambientes fortemente controlados, nem toda combinação entre
\textit{benchmark} e JVM necessariamente atinge um regime estável de
desempenho, o que torna arriscado assumir, \textit{a priori}, que poucas
iterações sejam suficientes para caracterizar o programa
\cite{vms-hot-and-cold}.

O conjunto de \textit{benchmarks} DaCapo Chopin foi escolhido por disponibilizar um conjunto diverso de cargas
de trabalho Java, o que reduz o risco de se tirar conclusões a partir de um
conjunto excessivamente restrito de aplicações \cite{dacapochopin}. Além disso, 
o pacote permite a inclusão de novas métricas por meio de \textit{callbacks} que são 
executadas antes e depois de cada iteração do \textit{benchmark}. Utilizando os 
\textit{callbacks}, foi possível coletar as métricas de energia e tempo de execução de cada iteração. 
Outra vantagem é que o DaCapo Chopin automaticamente executa as iterações de aquecimento e detecta 
falhas de convergência, o que permite descartar iterações que não representam o comportamento estável do \textit{benchmark}.

\section{Métricas de análise}

Além da energia do pacote e do tempo de execução, a análise considera duas
métricas derivadas: o \textit{Energy-Delay Product} (EDP), definido como
\(E \times T\), e o \(EDP^2\), definido como \(E^2 \times T\). Essas métricas
são importantes porque energia e tempo não variam necessariamente na mesma
direção. Uma configuração pode reduzir o tempo de execução, mas resultar em um
consumo energético maior. O EDP funciona como uma medida de equilíbrio entre
esses dois objetivos: ele penaliza tanto execuções lentas quanto execuções
energeticamente custosas, evitando que a análise considere apenas uma das
dimensões. Essa necessidade de observar energia e desempenho em conjunto já foi
apontada por \citeonline{programming_lang_energy}, que mostraram que
implementações mais rápidas não são necessariamente as mais eficientes em
termos energéticos. O \(EDP^2\) complementa o EDP ao aumentar o peso da
energia no produto. Isso penaliza mais as configurações com maior consumo
energético. Assim, quando duas configurações apresentam tempos próximos, mas
energias distintas, o \(EDP^2\) tende a evidenciar com mais clareza a vantagem
da alternativa energeticamente mais econômica. Métricas desse tipo são
utilizadas em outros contextos de avaliação de sistemas. O produto entre
energia e tempo é empregado, por exemplo, para orientar políticas de
escalonamento dinâmico de tensão e frequência \cite{green-governors}. Sendo
assim, o uso dessas métricas no atual contexto busca apresentar uma visão mais
consolidada da relação entre energia e tempo de execução e as configurações que
melhor otimizam esses dois objetivos.

\section{Trabalhos Relacionados}

A Tabela \ref{tab:trabalhos_relacionados} resume alguns trabalhos relacionados ao
tema e destaca, para cada caso, a diferença de escopo em relação a este trabalho.

\begin{table}[ht]
    \centering
    \caption{Comparação entre trabalhos relacionados e esta dissertação}
    \label{tab:trabalhos_relacionados}
    \small
    \renewcommand{\arraystretch}{1.2}
    \begin{tabularx}{\textwidth}{|p{2.5cm}|X|X|X|}
        \hline
        \textbf{Trabalho} & \textbf{Objeto de estudo} & \textbf{Métricas analisadas} & \textbf{Diferença principal em relação a este trabalho} \\
        \hline
        \citeonline{energy-java-collections-hasan} & Operações sobre coleções Java. & Energia. & Foca em microcenários de estruturas de dados, e não em cargas completas da JVM. \\
        \hline
        \citeonline{programming_lang_energy} & Problemas equivalentes em múltiplas linguagens. & Energia, tempo e memória. & Compara linguagens, enquanto esta dissertação analisa cargas Java e parâmetros da JVM. \\
        \hline
        \citeonline{jrapl} & Otimizações energéticas em programas Java e biblioteca JRAPL. & Energia e tempo. & Introduz a ferramenta de medição, mas não a integra a uma suíte consolidada como o DaCapo Chopin. \\
        \hline
        \citeonline{dalvik-vs-art} & Comparação entre Dalvik e ART em dispositivos móveis. & Energia e descarga de bateria. & Analisa \textit{runtimes} móveis, e não parâmetros de execução da JVM em cargas Java. \\
        \hline
        Este trabalho & Cargas do DaCapo Chopin na JVM com instrumentação energética integrada. & Energia do pacote, tempo, EDP e \(EDP^2\). & Integra a medição energética ao DaCapo Chopin e analisa como frequência e \textit{heap} afetam energia e desempenho. \\
        \hline
    \end{tabularx}
\end{table}

Esses trabalhos mostram que energia pode ser influenciada pela
escolha da linguagem, pela estrutura de dados utilizada, pelo \textit{runtime}
ou por otimizações específicas. Entretanto, permanece menos explorado como 
as diferentes frequências do processador e tamanhos de \textit{heap} interagem 
com diferentes cargas de trabalho e o impacto que isso gera sobre a energia e o tempo de execução. Esta
dissertação procura preencher essa lacuna ao combinar uma extensão do DaCapo
Chopin para coleta energética com uma análise comparativa de energia e tempo em
13 \textit{benchmarks} Java.
